%============================================================
% Appendix: Smoothness theory for discrete-continuous models
%============================================================
% Requires \usepackage{amsthm} in the preamble of dcegmSearch.tex.
% Input from Appendix.tex.
%============================================================

\theoremstyle{plain}
\newtheorem{lemma}{Lemma}[section]
\newtheorem{proposition}[lemma]{Proposition}
\newtheorem{corollary}[lemma]{Corollary}

\theoremstyle{definition}
\newtheorem{definition}[lemma]{Definition}
\newtheorem{assumption}{Assumption}[section]
\newtheorem{remark}[lemma]{Remark}

\newcommand{\keep}{\mathrm{keep}}
\newcommand{\adj}{\mathrm{adj}}

\section{A Theory of Smoothness in Discrete--Continuous Models}
\label{sec:appdx_smoothness}

%------------------------------------------------------------
\subsection{Purpose and roadmap}
\label{subsec:smooth_purpose}
%------------------------------------------------------------

The main text establishes by numerical illustration that taste shocks and directed
search reduce the discontinuities generated by the EGM step in discrete--continuous
(DC) models. This appendix replaces those illustrations with propositions. It does
three things: it shows \emph{why} the hard-max model fails
(Section~\ref{subsec:smooth_negative}), \emph{how} taste shocks repair it
(Section~\ref{subsec:smooth_taste}), and \emph{how} directed search repairs it
(Section~\ref{subsec:smooth_search}), in each case with the exact smoothness class,
an explicit bound on curvature, and an explicit bound on the model bias incurred.

Two preliminary remarks shape everything that follows.

First, ``smoothness'' is not the operative property. What determines whether the EGM
step produces suboptimal points is \emph{strict monotonicity of the endogenous grid
map} $a\mapsto m(a)$, and by Lemma~\ref{lem:egm} that is a condition on \emph{bounded
second derivatives} of the continuation value, not on differentiability as such.
Differentiability is a separate matter: it governs whether marginal propensities to
consume jump. Separating these two targets is what allows the claim in the abstract
(suboptimal points can be eliminated by search alone) and the claim in the conclusion
(uniform smoothness additionally requires taste shocks) to be stated as two distinct
and simultaneously correct propositions. Under directed search with the log-type cost
function~\eqref{eq:search_cost_basic}, the first holds and the second fails; see
Proposition~\ref{prop:search}.

Second, all three mechanisms are the same mathematical object evaluated at different
parameters. Each is the convex conjugate of a cost of randomizing over the discrete
alternatives (Lemma~\ref{lem:conjugate}). The smoothness class of the solved model is
governed entirely by the local behaviour of that cost at the boundary of the choice
simplex. This yields the sharpest result in the appendix,
Proposition~\ref{prop:tradeoff}: the apparent conflict between concentrating
adjustment near the $(S,s)$ thresholds and obtaining smooth decision rules is an
artifact of the particular cost function in~\eqref{eq:search_cost_basic}, not a
theorem, and both properties can be obtained simultaneously.

\begin{remark}[Notation]
\label{rem:notation}
Section~\ref{subsec:ModelBasicTaste} uses $\varepsilon_t^{\keep},\varepsilon_t^{\adj}$
for the Extreme Value Type~I draws while Section~\ref{subsec:ModelBasicSearch} uses
$\varepsilon_t$ for search intensity. To avoid the collision, this appendix reserves
$\varepsilon$ for search intensity and writes the Gumbel draws as
$\zeta_t^{\keep},\zeta_t^{\adj}$. The main text should be harmonised accordingly.
\end{remark}

%------------------------------------------------------------
\subsection{A common reduced form}
\label{subsec:smooth_framework}
%------------------------------------------------------------

Let $x\in X\subseteq\mathbb{R}^{n}$ denote the beginning-of-period continuous state:
$x=M_t$ in the retirement model of Section~\ref{sec:RetirementModel} and
$x=(m_t,n_t)$ in the durable model of Section~\ref{sec:ModelBasic}. Let
$j\in\{0,1\}$ index the two discrete branches, with \emph{conditional value functions}
$v_t^{0}(x)$ and $v_t^{1}(x)$, each the value of a continuous optimisation problem
solved conditional on $j$. In the retirement model $(v^0,v^1)=(V_t^{\Ret},V_t^{W})$;
in the durable model $(v^0,v^1)=(v_t^{\keep},v_t^{\adj}\circ\ell)$ with
$\ell(m,n)=m+(1-\tau)n$. Define the \emph{value gap} and the \emph{indifference set}
\begin{align}
\Delta_t(x) &\equiv v_t^{1}(x)-v_t^{0}(x),
&
\mathcal{K}_t &\equiv \{x\in X:\Delta_t(x)=0\}.
\end{align}

All three mechanisms resolve the discrete choice through an \emph{aggregator}
$\Lambda:\mathbb{R}^2\to\mathbb{R}$ applied pointwise,
$v_t(x)=\Lambda(v_t^{0}(x),v_t^{1}(x))$. Because each is translation equivariant,
$\Lambda(v^0+k,v^1+k)=\Lambda(v^0,v^1)+k$ for all $k\in\mathbb{R}$, every aggregator
admits the \emph{reduced form}
\begin{align}
v_t(x) &= v_t^{0}(x) + \Phi\big(\Delta_t(x)\big)
\label{eq:reduced_form}
\end{align}
for a single scalar function $\Phi:\mathbb{R}\to\mathbb{R}$ with $\Phi(0)$ finite. This
is the device that reduces every regularity question below to a statement about $\Phi$
near the origin, composed with a transversality condition on $\Delta_t$. The three
cases of interest are
\begin{align}
\Phi^{\max}(\delta) &= \max\{0,\delta\},
&
\Phi^{\sigma}(\delta) &= \sigma_\varepsilon\log\!\big(1+e^{\delta/\sigma_\varepsilon}\big),
&
\Phi^{S}(\delta) &= \max_{\varepsilon\in[0,1]}\big\{\varepsilon\delta-\sigma(\varepsilon)\big\},
\label{eq:three_Phi}
\end{align}
corresponding respectively to the hard maximum of Section~\ref{sec:ModelBasic}, the
Extreme Value Type~I log-sum~\eqref{eq:taste_logsum}, and the directed-search problem
\eqref{eq:search_discrete_basic} with cost~\eqref{eq:search_cost_basic}.

Because the hard-max case produces functions that are not globally differentiable, the
$C^{k}$ scale alone is insufficient bookkeeping. For a closed set
$\mathcal{N}\subset X$ of measure zero, let $\mathcal{PC}^{k}(\mathcal{N})$ denote the
functions that are $C^{k}$ on $X\setminus\mathcal{N}$ with finite one-sided limits of
all derivatives up to order $k$ at each point of $\mathcal{N}$. Write
$\mathcal{N}_t$ for the \emph{singular set} of $v_t$, i.e.\ the set at which $v_t$
fails to be $C^1$. Tracking the cardinality (or the surface measure) of
$\mathcal{N}_t$ across backward iterations is what formalises the statement in
Section~\ref{sec:RetirementModel} that kinks are ``propagated back in time, thus,
exacerbating the problem.''

Finally, let
\begin{align}
\mathcal{E}_t(a) &\equiv \E_t\big[v_{m,t+1}(\mathcal{R}a+\theta_{t+1})\big],
&
m_t(a) &\equiv a + c_t(a),
\end{align}
so that $\mathcal{E}_t$ is the object the EGM step inverts and $m_t$ is the endogenous
grid map. The following definition makes the paper's central claim precise.

\begin{definition}
\label{def:clean}
The EGM step at $t$ is \emph{clean} if $a\mapsto m_t(a)$ is continuous and strictly
increasing on the unconstrained region. If the EGM step is clean, the candidate points
returned by the Euler equation require no upper-envelope refinement.
\end{definition}

%------------------------------------------------------------
\subsection{Assumptions}
\label{subsec:smooth_assumptions}
%------------------------------------------------------------

\begin{assumption}[Felicity]
\label{as:u}
$u$ is $C^{\infty}$ on the interior of its domain, strictly increasing and strictly
concave in each argument, with Inada conditions in each argument. In the durable
model $u(c,d)=(c^{\alpha}d^{1-\alpha})^{1-\rho}/(1-\rho)$ with
$\alpha\in(0,1)$, $\rho>0$; this $u$ is concave, being the composition of the
increasing concave map $z\mapsto z^{1-\rho}/(1-\rho)$ with the concave,
degree-one homogeneous aggregator $c^{\alpha}d^{1-\alpha}$.
\end{assumption}

\begin{assumption}[Interiority and regular constraint set]
\label{as:interior}
For each $t$ and each $x$ in the region of interest the optimal continuous choice is
interior, and the set on which $a_t\ge 0$ binds is relatively open with a $C^1$
boundary. Moreover the conditional marginal propensity to consume satisfies
$0<\partial c_t^{j}/\partial m<1$ on the unconstrained region.
\end{assumption}

Assumption~\ref{as:interior} is not cosmetic. The borrowing constraint generates a
kink of its own at the constrained/unconstrained boundary, and any statement that
kinks in the reported policy functions are \emph{discrete-choice} kinks requires that
this second source be separated out. The bound $\partial c/\partial m<1$ is used only
in Proposition~\ref{prop:uniform}.

\begin{assumption}[Terminal regularity]
\label{as:terminal}
At $T$ both conditional value functions are $C^{\infty}$ on the interior of the state
space.
\end{assumption}

Assumption~\ref{as:terminal} holds in both applications and supplies the base case of
every induction below: the retirement terminal problem is the closed-form logarithmic
problem, and the durable terminal problem has $c_T^{\keep}=m_T$, $c_T^{\adj}=\alpha\ell_T$
and $d_T^{\adj}=(1-\alpha)\ell_T/p_T$.

\begin{assumption}[Transversality]
\label{as:transversal}
For each $t$, $\mathcal{K}_t$ is a nonempty $C^{1}$ embedded submanifold of $X$ of
codimension one, and $\nabla\Delta_t(x)\neq 0$ for every $x\in\mathcal{K}_t$. Define
the \emph{jump magnitude}
\begin{align}
\kappa_t &\equiv \inf_{x\in\mathcal{K}_t}\big|\nabla\Delta_t(x)\big| > 0,
&
\Xi_t &\equiv \sup_{x\in\mathcal{K}_t}\big|\nabla\Delta_t(x)\big| < \infty .
\end{align}
In the retirement model this reduces to $\lambda_t^{W}(M^{\ast})\neq\lambda_t^{\Ret}(M^{\ast})$
at each crossing point $M^{\ast}$; in the durable model
$\nabla\Delta_t=\big(\lambda_t^{\adj}-\lambda_t^{\keep},\,(1-\tau)\lambda_t^{\adj}-v_{n,t}^{\keep}\big)$.
\end{assumption}

Assumption~\ref{as:transversal} does double duty. It is what makes the negative result
\emph{true} --- a degenerate tangential crossing would produce no kink at all --- and
it is what guarantees that the smoothing region in
Propositions~\ref{prop:evtaste} and~\ref{prop:search} has positive width of order
$\sigma_\varepsilon$ and $S$ respectively.

\begin{assumption}[Single crossing]
\label{as:singlecross}
For each fixed $n$, $\Delta_t(\cdot,n)$ has finitely many zeros, and in the durable
model $\partial_m\Delta_t>0$ on $\mathcal{K}_t$.
\end{assumption}

Assumption~\ref{as:singlecross} is what delivers an $(S,s)$ \emph{band} rather than an
arbitrary adjustment set. It is the least innocent assumption in this list: proving it
from primitives in a model whose continuation value is itself an aggregate of two
branches is not attempted here, and it should be verified numerically on the solved
model (Section~\ref{subsec:smooth_tests}).

\begin{assumption}[Shock regularity]
\label{as:shocks}
The transitory income shock $\theta_{t}$ has a density $f_\theta\in C^{\infty}$ with
all derivatives integrable against the kernels appearing below; the mean-one
log-normal specification of Sections~\ref{sec:RetirementModel}
and~\ref{sec:ModelBasic} qualifies. Where a taste shock $\xi_t$ is present its
density $f_\xi$ satisfies the same conditions.
\end{assumption}

Assumption~\ref{as:shocks} refers to the \emph{model}. The \emph{computed}
approximation replaces these densities by $J$- and $K$-node quadratures.
Lemma~\ref{lem:mollify} shows that these two objects have entirely different
regularity, which is the rigorous content of the phantom-kink discussion in
Section~\ref{sec:incdisc}.

\begin{assumption}[Impatience]
\label{as:impatience}
$\beta\mathcal{R}^{2}<1$.
\end{assumption}

Assumption~\ref{as:impatience} is required only for the uniform-in-$t$ estimate of
Proposition~\ref{prop:uniform} and is satisfied by both calibrations used in the
paper: $\mathcal{R}=1,\beta=0.98$ gives $\beta\mathcal{R}^2=0.98$, and
$\mathcal{R}=1.04^{0.25},\beta=0.975$ gives $\beta\mathcal{R}^2\approx 0.994$.

\begin{assumption}[Level normalisation; retirement model only]
\label{as:normalisation}
The felicity function is normalised so that the multiplicative taste shock
$\xi_t V_t^{W}$ of Section~\ref{sec:RetirementModel} is well defined, i.e.\ the
additive constant in $u$ is fixed.
\end{assumption}

Assumption~\ref{as:normalisation} deserves emphasis because it is substantive rather
than technical. A value function is determined only up to an affine transformation of
felicity, so a shock entering multiplicatively on $V_t^{W}$ is \emph{not} invariant to
that transformation: adding a constant to $u$ changes the model's predictions. The
accompanying code fixes the normalisation with an additive constant chosen for
interpolation accuracy. Any proposition about the multiplicative log-normal taste
shock is therefore conditional on Assumption~\ref{as:normalisation}. The Extreme Value
Type~I specification of Section~\ref{subsec:ModelBasicTaste}, being additive, requires
no such assumption, which is an argument in its favour that the main text does not
currently make.

%------------------------------------------------------------
\subsection{Four lemmas}
\label{subsec:smooth_lemmas}
%------------------------------------------------------------

\begin{lemma}[Regularity transfer and the envelope bootstrap]
\label{lem:transfer}
Under Assumptions~\ref{as:u}--\ref{as:interior}, for any $k\ge 0$:
\begin{enumerate}
    \item $\mathcal{E}_t\in C^{k}\;\Longrightarrow\;c_t=(u')^{-1}(\beta\mathcal{R}\mathcal{E}_t)\in C^{k}$;
    \item $c_t^{j}\in C^{k}\;\Longrightarrow\;v_t^{j}\in C^{k+1}$;
    \item consequently $v_{t+1}\in C^{k+1}\;\Longrightarrow\;v_t^{j}\in C^{k+1}$ for $j=0,1$.
\end{enumerate}
\end{lemma}

\begin{proof}
(i) $u'$ is $C^{\infty}$ with $u''<0$ on the interior, so by the inverse function
theorem $(u')^{-1}$ is $C^{\infty}$ on its range; a composition of a $C^{\infty}$ map
with a $C^{k}$ map is $C^{k}$. (ii) The envelope condition gives
$v_{m,t}^{j}=u'(c_t^{j})$, which is $C^{k}$ by the same argument; a function whose
gradient is $C^{k}$ is $C^{k+1}$. (iii) $v_{t+1}\in C^{k+1}$ implies
$v_{m,t+1}\in C^{k}$, hence $\mathcal{E}_t\in C^{k}$ (a finite sum or an integral of
$C^{k}$ functions), hence $c_t\in C^{k}$ by (i), hence $v_t^{j}\in C^{k+1}$ by (ii).
\end{proof}

Lemma~\ref{lem:transfer}(iii) is the workhorse of the appendix: the \emph{envelope
condition buys back the derivative lost in the Euler equation}, so conditional value
functions preserve their smoothness class across a backward step. It follows that
\emph{only the aggregator can destroy smoothness}, and the entire induction reduces to
the question of what class $\Phi$ admits.

\begin{lemma}[Mollification: exact integration versus quadrature]
\label{lem:mollify}
Let $g$ be locally bounded and piecewise $C^{\infty}$ with finitely many jump
discontinuities in $g'$, of magnitudes $[g'](z_i)$ at points $z_i$. Fix $y>0$.
\begin{enumerate}
    \item \emph{(Exact integration smooths infinitely.)} Under Assumption~\ref{as:shocks},
    $G(b)\equiv\int g(b+y\eta)f(\eta)\,d\eta\in C^{\infty}$, regardless of the
    singularities of $g$.
    \item \emph{(Quadrature does not smooth at all.)} For
    $G_J(b)\equiv\sum_{j=1}^{J}\omega_j\,g(b+y\eta_j)$ with $\omega_j>0$,
    $\sum_j\omega_j=1$, the derivative $G_J'$ has a jump of magnitude
    $\omega_j\,[g'](z_i)$ at each $b$ with $b+y\eta_j=z_i$. Hence $G_J\notin C^{1}$
    for every $J$.
    \item \emph{(Conservation of total jump.)} If the preimages are distinct, then for
    every $J$,
    \begin{align}
    \sum_{i}\sum_{j}\big|[G_J']\big| \;=\; \sum_{i}\big|[g'](z_i)\big| ,
    \end{align}
    so the \emph{largest} jump falls at rate $O(\max_j\omega_j)=O(1/J)$ while the
    \emph{total} jump is invariant in $J$.
\end{enumerate}
\end{lemma}

\begin{proof}
(i) Substituting $z=b+y\eta$ gives $G(b)=y^{-1}\int g(z)f\big((z-b)/y\big)dz$.
Differentiation in $b$ falls on $f$, not on $g$:
$G^{(r)}(b)=(-1)^{r}y^{-(r+1)}\int g(z)f^{(r)}\big((z-b)/y\big)dz$, which is finite for
every $r$ under Assumption~\ref{as:shocks} and continuous in $b$ by dominated
convergence. (ii) Immediate from $G_J'(b)=\sum_j\omega_j g'(b+y\eta_j)$ wherever
defined, evaluating one-sided limits. (iii) Sum the jumps in (ii) over $j$ and use
$\sum_j\omega_j=1$.
\end{proof}

Lemma~\ref{lem:mollify} is the precise version of the phantom-kink claim of
Section~\ref{sec:incdisc}. Part (iii) says that refining the income grid makes the
solution \emph{look} smooth without ever making it smooth, and it predicts exactly the
pattern in Table~\ref{tab:incdisc}: the mean absolute deviation collapses in $J$ while
the maximum absolute deviation falls slowly and non-monotonically. It also supplies a
better diagnostic than a maximum absolute deviation, namely the total variation of
$c_t'$, which is invariant in $J$ under quadrature and finite under exact integration.

\begin{lemma}[EGM monotonicity criterion]
\label{lem:egm}
Wherever $c_t$ is differentiable,
\begin{align}
m_t'(a) \;=\; 1+c_t'(a) \;=\; 1 - \frac{\beta\mathcal{R}^{2}\,\E_t\big[v_{mm,t+1}\big]}{\big|u''(c_t)\big|} .
\label{eq:egm_slope}
\end{align}
Consequently:
\begin{enumerate}
    \item if $\E_t[v_{mm,t+1}]\le 0$ then $m_t'(a)\ge 1$, and the EGM step is clean;
    \item in general the EGM step is clean if and only if
    \begin{align}
    \beta\mathcal{R}^{2}\,\E_t\big[v_{mm,t+1}\big] \;<\; \big|u''(c_t)\big|
    \qquad\text{pointwise;}
    \label{eq:egm_condition}
    \end{align}
    \item if $\mathcal{E}_t$ has a jump discontinuity then $m_t$ has a jump and the EGM
    step is not clean.
\end{enumerate}
\end{lemma}

\begin{proof}
Differentiate $u'(c_t(a))=\beta\mathcal{R}\mathcal{E}_t(a)$ in $a$, using
$\mathcal{E}_t'(a)=\mathcal{R}\,\E_t[v_{mm,t+1}]$, to obtain
$c_t'(a)=\beta\mathcal{R}^{2}\E_t[v_{mm,t+1}]/u''(c_t)$; substitute $u''=-|u''|$ and add
one. Parts (i)--(iii) follow by inspection of~\eqref{eq:egm_slope}.
\end{proof}

Condition~\eqref{eq:egm_condition} is the operative criterion of this appendix. It
converts ``the solution is smooth enough'' into a checkable inequality, and combined
with Proposition~\ref{prop:uniform} it yields an explicit threshold on the smoothing
parameter. Two caveats on how to state it usefully. First, under CRRA felicity
$\inf|u''|$ vanishes at high wealth, so~\eqref{eq:egm_condition} is vacuous if stated
globally; it should be stated on a compact wealth region, or, preferably and in
keeping with the implementation, in inverse-marginal-utility
(consumption-equivalent) units so that both sides are scale free. Second,
Lemma~\ref{lem:egm} presumes differentiability of $c_t$; where $c_t$ merely has a kink,
$m_t'$ is to be read as the pair of one-sided derivatives and the criterion applied to
each.

\begin{lemma}[Aggregator regularity via conjugate duality]
\label{lem:conjugate}
Let $\sigma:[0,1]\to\mathbb{R}$ be strictly convex, $C^{2}$ on $(0,1)$, with
$\sigma(0)=0$, and let $\Phi=\sigma^{\ast}$ be its convex conjugate as
in~\eqref{eq:three_Phi}. Then:
\begin{enumerate}
    \item $\Phi$ is convex and nondecreasing with $\Phi'(\delta)=\varepsilon^{\ast}(\delta)\in[0,1]$,
    and $\varepsilon^{\ast}=(\sigma')^{-1}(\delta)$ wherever interior;
    \item $\varepsilon^{\ast}(\delta)=0$ for all $\delta\le 0$ if and only if $\sigma'(0)=0$
    (\emph{one-sidedness});
    \item given $\sigma'(0)=0$, $\;\Phi''(0^{+})=1/\sigma''(0^{+})$ and
    $\Phi''(0^{-})=0$, so
    \begin{align}
    \Phi\in C^{2}\quad\Longleftrightarrow\quad \sigma''(0^{+})=+\infty ;
    \end{align}
    \item if $\sigma'(\varepsilon)= A\varepsilon^{1/\gamma}$ near $\varepsilon=0$ for an
    integer $\gamma\ge 1$ and $A>0$, then $\varepsilon^{\ast}(\delta)=(\delta^{+}/A)^{\gamma}$
    near $\delta=0$ and $\Phi\in C^{\gamma}\setminus C^{\gamma+1}$.
\end{enumerate}
\end{lemma}

\begin{proof}
(i) A conjugate is a supremum of affine functions of $\delta$, hence convex;
monotonicity follows from $\varepsilon\ge0$; $\Phi'=\varepsilon^{\ast}$ is Danskin's
theorem, and the interior first-order condition is $\delta=\sigma'(\varepsilon^{\ast})$.
(ii) The constraint $\varepsilon\ge0$ binds exactly when $\delta\le\sigma'(0)$.
(iii) Differentiating $\delta=\sigma'(\varepsilon^{\ast}(\delta))$ gives
$\Phi''(\delta)=1/\sigma''(\varepsilon^{\ast}(\delta))$, and $\varepsilon^{\ast}(0^{+})=0$;
on $\delta<0$, $\Phi\equiv0$. (iv) Invert $\sigma'$ and note that
$\delta\mapsto(\delta^{+})^{\gamma}$ is $C^{\gamma-1}$ but not $C^{\gamma}$ at the
origin, so $\Phi'\in C^{\gamma-1}\setminus C^{\gamma}$.
\end{proof}

Lemma~\ref{lem:conjugate} places both smoothing devices in a single family. The
Extreme Value Type~I log-sum is the conjugate of the Shannon entropy on the two-point
simplex, and directed search is the conjugate of the search cost. This is the
additive-perturbed-utility representation of stochastic discrete choice studied by
Mattsson and Weibull \citeyearpar{Mattsson2002-pc}, Hofbauer and Sandholm
\citeyearpar{Hofbauer2002-gc} and \citet{Fudenberg2015-sv}; see
\citet{Rockafellar1970-cv} for the duality facts used above. Two implications matter
for the paper. First, ``rationalised'' versus ``behavioural'' is not a structural
distinction between the two mechanisms, since both are conjugates of a control cost
and the associated random-utility and perturbed-utility representations are
observationally equivalent; the durable arguments for search should therefore rest on
one-sidedness (Proposition~\ref{prop:onesided}) and on testability, which do survive.
Second, and constructively, the modeller \emph{chooses} $\sigma$, so by
Lemma~\ref{lem:conjugate}(iii)--(iv) the modeller chooses the smoothness class; see
Proposition~\ref{prop:tradeoff}.

%------------------------------------------------------------
\subsection{Why the hard-max model is not smooth}
\label{subsec:smooth_negative}
%------------------------------------------------------------

\begin{proposition}[Convex kink and marginal-value jump]
\label{prop:kink}
Let $\Phi=\Phi^{\max}$ and suppose Assumptions~\ref{as:u}--\ref{as:transversal} hold
with $v_t^{0},v_t^{1}\in C^{1}$. Then $v_t$ is Lipschitz, $v_t\notin C^{1}$ at every
$x\in\mathcal{K}_t$, and along any $C^{1}$ curve $\gamma$ crossing $\mathcal{K}_t$
transversally at $\gamma(0)=x^{\ast}$ in the direction of increasing $\Delta_t$,
\begin{align}
\frac{d}{ds}v_t(\gamma(s))\Big|_{s=0^{+}} - \frac{d}{ds}v_t(\gamma(s))\Big|_{s=0^{-}}
\;=\; \frac{d}{ds}\Delta_t(\gamma(s))\Big|_{s=0} \;>\;0 .
\label{eq:kink_jump}
\end{align}
The jump is therefore strictly \emph{upward}: $v_t$ has a convex kink and, in the
distributional sense, $v_{mm,t}$ carries a strictly positive Dirac mass on
$\mathcal{K}_t$.
\end{proposition}

\begin{proof}
Write $h(s)=v_t(\gamma(s))=v_t^{0}(\gamma(s))+\max\{0,\Delta_t(\gamma(s))\}$. For
$s<0$, $\Delta_t(\gamma(s))<0$ and $h'(s)=\tfrac{d}{ds}v_t^{0}$; for $s>0$,
$h'(s)=\tfrac{d}{ds}v_t^{0}+\tfrac{d}{ds}\Delta_t=\tfrac{d}{ds}v_t^{1}$. Taking
one-sided limits at $s=0$ and differencing gives~\eqref{eq:kink_jump}, which is
strictly positive by transversality (Assumption~\ref{as:transversal}) and the choice of
direction. Lipschitz continuity is inherited from the branches.
\end{proof}

Proposition~\ref{prop:kink} identifies the jump magnitude with $|\nabla\Delta_t|$,
i.e.\ with the difference in the branches' marginal values,
$\kappa_t=|\lambda^{1}-\lambda^{0}|$ in the retirement model. The \emph{sign} matters:
because the maximum always selects the steeper branch after a crossing, the kink is
convex and $\E_t[v_{mm,t+1}]$ acquires a positive singular part, which is precisely the
direction that violates~\eqref{eq:egm_condition}.

\begin{proposition}[Quadrature guarantees failure]
\label{prop:quadfail}
Let $\Phi=\Phi^{\max}$ and let the income expectation be computed by $J$-node
quadrature. Then $\mathcal{E}_{t-1}$ has a jump of magnitude $\omega_j\kappa_t$ at each
$a$ with $\mathcal{R}a+y\eta_j\in\mathcal{K}_t$. Consequently $c_{t-1}$ has at least
$J\cdot|\mathcal{K}_t|$ jump discontinuities, $m_{t-1}$ is discontinuous, and by
Lemma~\ref{lem:egm}(iii) the EGM step is not clean --- for every $J$ and every
parameterisation satisfying Assumption~\ref{as:transversal}.
\end{proposition}

\begin{proof}
Apply Lemma~\ref{lem:mollify}(ii) with $g=v_{m,t}$, whose jumps are given by
Proposition~\ref{prop:kink}, then Lemma~\ref{lem:transfer}(i) and
Lemma~\ref{lem:egm}(iii).
\end{proof}

\begin{proposition}[Accumulation]
\label{prop:accumulate}
Under the hypotheses of Proposition~\ref{prop:quadfail}, the singular set satisfies
\begin{align}
\mathcal{N}_{t-1} \;\supseteq\; \mathcal{K}_{t-1}\;\cup\;
\big\{x : \mathcal{R}a(x)+y\eta_j \in \mathcal{K}_{t}\ \text{for some } j\le J\big\},
\end{align}
so $|\mathcal{N}_{t-1}|\ge |\mathcal{K}_{t-1}| + J\,|\mathcal{K}_{t}|$ and the number of
points at which the conditional value functions fail to be $C^{1}$ is nondecreasing,
and generically strictly increasing, as $t$ recedes from $T$.
\end{proposition}

\begin{proof}
By Proposition~\ref{prop:quadfail} the conditional policy $c_{t-1}^{j}$ has jumps at
the stated points, so by the envelope condition $v_{m,t-1}^{j}=u'(c_{t-1}^{j})$ has
jumps there and $v_{t-1}^{j}\notin C^{1}$; the new indifference set
$\mathcal{K}_{t-1}$ contributes additionally by Proposition~\ref{prop:kink}. Note that
Lemma~\ref{lem:transfer}(iii) does not apply, since it requires $k\ge0$ and here
$v_{m,t}\notin C^{0}$.
\end{proof}

Proposition~\ref{prop:accumulate} is the formal content of the informal claim that
kinks propagate backwards and worsen the problem. It also shows precisely where
Lemma~\ref{lem:transfer} breaks: the envelope bootstrap preserves smoothness only from
$C^{0}$ upwards, and the hard maximum pushes the solution below that floor.

\begin{proposition}[Exact income integration is necessary but not sufficient]
\label{prop:exactincome}
Let $\Phi=\Phi^{\max}$ and let the income expectation be exact under
Assumption~\ref{as:shocks}. Then $\mathcal{E}_{t-1}\in C^{\infty}$ and $c_{t-1}$ is
continuous, but
\begin{align}
\E_{t-1}\big[v_{mm,t}\big] \;\ni\; \frac{\kappa_t}{y}\,f_\theta(\eta^{\ast}),
\qquad \eta^{\ast}=\frac{M^{\ast}-\mathcal{R}a}{y},
\end{align}
as a spike at each preimage of $\mathcal{K}_t$, so by~\eqref{eq:egm_condition} the EGM
step fails to be clean whenever
\begin{align}
\beta\mathcal{R}^{2}\,\frac{\kappa_t}{y}\,f_\theta(\eta^{\ast})
\;\ge\; \big|u''(c_{t-1})\big| .
\label{eq:income_threshold}
\end{align}
\end{proposition}

\begin{proof}
Continuity and smoothness of $\mathcal{E}_{t-1}$ follow from
Lemma~\ref{lem:mollify}(i). Decompose $v_{m,t}=\tilde v_{m,t}+\kappa_t\mathbbm{1}\{z\ge M^{\ast}\}$
with $\tilde v_{m,t}$ continuous. The indicator contributes
$\kappa_t\Pr[\mathcal{R}a+y\theta\ge M^{\ast}]$ to $\mathcal{E}_{t-1}$, whose derivative in
$a$ is $\kappa_t\mathcal{R}f_\theta(\eta^{\ast})/y$; dividing by $\mathcal{R}$ gives the
stated contribution to $\E[v_{mm,t}]$. Apply Lemma~\ref{lem:egm}(ii).
\end{proof}

Proposition~\ref{prop:exactincome} rationalises the finding in
Figs.~\ref{fig:fig4}(a)--(b) that income uncertainty adds smoothness but does not
remove the discontinuities: exact integration upgrades the policy from discontinuous to
continuous, but leaves a curvature spike whose height is inversely proportional to the
dispersion of income. The discrete choice must be smoothed as well.

\begin{remark}[Two distinct discontinuities]
\label{rem:twodiscont}
Two different objects are non-smooth under a hard maximum and the main text does not
consistently distinguish them.
\begin{enumerate}
    \item The \emph{combined} policy $c_t=(1-\mathbbm{1}_t)c_t^{\keep}+\mathbbm{1}_t c_t^{\adj}$
    jumps on $\mathcal{K}_t$ because the argmax switches branch. This is genuine
    economics; no solution method removes it, and a smoothing device removes it only by
    reporting an ex-ante expected policy in its place.
    \item The \emph{conditional} policy $c_t^{j}$ jumps because $v_{m,t+1}$ jumps
    (Proposition~\ref{prop:quadfail}). This is what breaks the EGM step.
\end{enumerate}
The figures in Sections~\ref{sec:RetirementModel} and~\ref{sec:DurableModel} plot
objects of type (ii), and the text should say so explicitly; otherwise the paper is
open to the objection that it smooths away the economics rather than the numerics.
\end{remark}

%------------------------------------------------------------
\subsection{Taste shocks}
\label{subsec:smooth_taste}
%------------------------------------------------------------

\subsubsection{Extreme Value Type~I}

\begin{proposition}[EV Type~I delivers $C^{\infty}$]
\label{prop:evtaste}
Let $\Phi=\Phi^{\sigma}$ with scale $\sigma_\varepsilon>0$. Then $\Phi^{\sigma}$ is real
analytic on $\mathbb{R}$, with
\begin{align}
(\Phi^{\sigma})'(\delta) = P^{\adj}=\frac{1}{1+e^{-\delta/\sigma_\varepsilon}}\in(0,1),
\qquad
(\Phi^{\sigma})''(\delta) = \frac{P^{\adj}(1-P^{\adj})}{\sigma_\varepsilon}\le\frac{1}{4\sigma_\varepsilon},
\end{align}
and $0<\Phi^{\sigma}(\delta)-\Phi^{\max}(\delta)=\sigma_\varepsilon\log\big(1+e^{-|\delta|}\big)\le\sigma_\varepsilon\log 2$.
Under Assumptions~\ref{as:u}--\ref{as:shocks} and exact income integration,
$v_t\in C^{\infty}$ for every $t\le T$, $c_t\in C^{\infty}$, marginal propensities to
consume are continuous, and
\begin{align}
\big\|v_{mm,t}\big\|_{\infty} \;\le\; \max_{j}\big\|v_{mm,t}^{j}\big\|_{\infty}
+ \frac{\Xi_t^{2}}{4\sigma_\varepsilon}.
\label{eq:ev_curvature}
\end{align}
\end{proposition}

\begin{proof}
The derivative formulas are direct, and $P(1-P)\le\tfrac14$. The bias identity follows
from $\sigma_\varepsilon\log(1+e^{\delta/\sigma_\varepsilon})-\max\{0,\delta\}
=\sigma_\varepsilon\log(1+e^{-|\delta|})$. For the class: by
Assumption~\ref{as:terminal} the base case holds at $T$; if $v_{t+1}\in C^{\infty}$
then $v_t^{j}\in C^{\infty}$ by Lemma~\ref{lem:transfer}(iii), and
$v_t=v_t^{0}+\Phi^{\sigma}(\Delta_t)$ is a composition of real analytic with
$C^{\infty}$ maps. Differentiating~\eqref{eq:reduced_form} twice gives
$v_{mm,t}=(1-P)v_{mm,t}^{0}+Pv_{mm,t}^{1}+(\Phi^{\sigma})''(\Delta_t)(\partial_m\Delta_t)^{2}$,
and bounding each term yields~\eqref{eq:ev_curvature}.
\end{proof}

Proposition~\ref{prop:evtaste} combined with Lemma~\ref{lem:egm} yields the explicit
\emph{bias--smoothness trade-off} that the main text asserts only qualitatively:
curvature is $O(1/\sigma_\varepsilon)$ while the per-period model bias is
$O(\sigma_\varepsilon)$. Solving~\eqref{eq:egm_condition} for $\sigma_\varepsilon$
delivers a threshold $\sigma_\varepsilon^{\ast}$ above which the EGM step is clean, and
$\sigma_\varepsilon^{\ast}\log 2$ is then the model distortion required to purchase an
envelope-free solve. That number is the honest headline statistic for the exercise and
should be reported.

\subsubsection{Multiplicative log-normal}

\begin{proposition}[Discretised log-normal taste shocks do not smooth]
\label{prop:lognormaltaste}
Consider the retirement specification of Section~\ref{sec:RetirementModel} under
Assumption~\ref{as:normalisation}.
\begin{enumerate}
    \item With \emph{exact} integration over $\xi$,
    $\E_\xi\big[\max\{V_{t+1}^{\Ret},\xi V_{t+1}^{W}\}\big]\in C^{\infty}$ by
    Lemma~\ref{lem:mollify}(i).
    \item With $K$-node discretisation, the object actually computed,
    \begin{align}
    \goth{V}_t^{1}(A)=\sum_{j=1}^{J}\sum_{k=1}^{K}\omega_j\mu_k
    \max\big\{V_{t+1}^{\Ret}(M^{j}),\,\xi_k V_{t+1}^{W}(M^{j})\big\},
    \end{align}
    is Lipschitz but not $C^{1}$: it has $J\!\cdot\!K\!\cdot\!|\mathcal{K}_{t+1}|$ kinks,
    of individual magnitude $\omega_j\mu_k\kappa$, and by
    Lemma~\ref{lem:mollify}(iii) the total jump is invariant in $K$.
\end{enumerate}
\end{proposition}

Proposition~\ref{prop:lognormaltaste} is a negative result for the paper as currently
written. Section~\ref{sec:RetirementModel} presents the ability of the log-normal shock
to ``take up multiple distinct values in the discrete approximation'' as an advantage
over Extreme Value Type~I shocks. For the purpose of smoothing it is a
\emph{disadvantage}: the EV Type~I assumption integrates out in closed form and yields
an exactly $C^{\infty}$ aggregator with no quadrature at all, whereas a $K$-point
log-normal yields an approximation that is never $C^{1}$ and whose total roughness does
not decline with $K$. Only the largest individual kink declines, at rate $O(1/K)$,
which is why very large taste-shock node counts are required in practice. Together with
the non-invariance flagged in Assumption~\ref{as:normalisation}, this argues for
replacing the multiplicative log-normal specification by an additive one.

%------------------------------------------------------------
\subsection{Directed search}
\label{subsec:smooth_search}
%------------------------------------------------------------

\begin{lemma}[Closed form of the search aggregator]
\label{lem:searchclosed}
For $\sigma(\varepsilon)=S[\varepsilon+(1-\varepsilon)\log(1-\varepsilon)]$ as
in~\eqref{eq:search_cost_basic}, the conjugate is
\begin{align}
\Phi^{S}(\delta)=
\begin{cases}
\delta - S\big(1-e^{-\delta/S}\big), & \delta>0,\\[2pt]
0, & \delta\le 0,
\end{cases}
\qquad
\varepsilon^{\ast}(\delta)=\big(\Phi^{S}\big)'(\delta)=\max\big\{0,\,1-e^{-\delta/S}\big\},
\label{eq:PhiS}
\end{align}
and
\begin{align}
\Phi^{S}(0)=0,\quad
\big(\Phi^{S}\big)'(0^{\pm})=0,\quad
\big(\Phi^{S}\big)''(0^{-})=0,\quad
\big(\Phi^{S}\big)''(0^{+})=\frac{1}{S},\quad
-S<\Phi^{S}-\Phi^{\max}\le 0 .
\end{align}
\end{lemma}

\begin{proof}
The interior first-order condition $\delta=\sigma'(\varepsilon)=-S\log(1-\varepsilon)$
gives $1-\varepsilon^{\ast}=e^{-\delta/S}$. Substituting,
$\sigma(\varepsilon^{\ast})=S\varepsilon^{\ast}-\delta e^{-\delta/S}$, hence
$\Phi^{S}=\varepsilon^{\ast}\delta-\sigma(\varepsilon^{\ast})
=\delta(\varepsilon^{\ast}+e^{-\delta/S})-S\varepsilon^{\ast}=\delta-S(1-e^{-\delta/S})$.
The derivative values follow from $(\Phi^{S})''(\delta)=S^{-1}e^{-\delta/S}$ on
$\delta>0$; the bias bound from $\Phi^{S}(\delta)-\delta=-S(1-e^{-\delta/S})\in(-S,0)$
for $\delta>0$ and equality to zero for $\delta\le0$.
\end{proof}

Note the contrast in the direction of the bias. The log-sum lies \emph{above} the
maximum, by at most $\sigma_\varepsilon\log 2$, reflecting the option value of taste
variation; directed search lies \emph{below} it, by at most $S$, reflecting the
resources absorbed by search. Both biases vanish linearly in the smoothing parameter.

\begin{proposition}[Directed search delivers exactly $C^{1}$, invariantly in $t$]
\label{prop:search}
Let $\Phi=\Phi^{S}$ with $S>0$ and suppose
Assumptions~\ref{as:u}--\ref{as:shocks} hold with exact income integration. Then for
every $t\le T$:
\begin{enumerate}
    \item $v_t\in C^{1}$ and, on $\mathcal{K}_t$, $v_t\notin C^{2}$: the class $C^{1}$
    is \emph{invariant} under backward iteration;
    \item $v_{m,t}$ is continuous, hence $\mathcal{E}_{t-1}$ is continuous and
    $m_{t-1}$ is continuous: the failure mechanism of
    Proposition~\ref{prop:quadfail} is eliminated, and the EGM step is clean whenever
    \eqref{eq:egm_condition} holds;
    \item $c_{t-1}\in C^{0}\setminus C^{1}$: the conditional policy is continuous but
    marginal propensities to consume jump, by
    $\big(\partial_m\Delta_t\big)^{2}/S$ across $\mathcal{K}_t$;
    \item $\|(\Phi^{S})''\|_{\infty}=1/S$, so
    $\|v_{mm,t}\|_{\infty}\le\max_j\|v_{mm,t}^{j}\|_{\infty}+\Xi_t^{2}/S$.
\end{enumerate}
\end{proposition}

\begin{proof}
(i) Induction. The base case at $T$ is Assumption~\ref{as:terminal}. Suppose
$v_{t+1}\in C^{1}$. Then $v_t^{j}\in C^{1}$ by Lemma~\ref{lem:transfer}(iii) with
$k=0$, and in fact $v_t^{j}\in C^{\infty}$ on the complement of the inherited singular
set. By Lemma~\ref{lem:searchclosed}, $\Phi^{S}\in C^{1}(\mathbb{R})$, so
$v_t=v_t^{0}+\Phi^{S}(\Delta_t)\in C^{1}$; and $\min\{k,1\}=1$, so the class does not
degrade further. Failure of $C^{2}$ on $\mathcal{K}_t$ follows from
$(\Phi^{S})''(0^{+})-(\Phi^{S})''(0^{-})=1/S\neq0$ together with
Assumption~\ref{as:transversal}. (ii) Immediate from (i) and
Lemma~\ref{lem:transfer}(i). (iii) $c_{t-1}=(u')^{-1}(\beta\mathcal{R}\mathcal{E}_{t-1})$
inherits the class of $\mathcal{E}_{t-1}$, which is that of $v_{m,t}$, namely
$C^{0}\setminus C^{1}$; the size of the MPC jump follows from differentiating
$v_{mm,t}$ twice through~\eqref{eq:reduced_form}. (iv) From
Lemma~\ref{lem:searchclosed} and the second derivative of~\eqref{eq:reduced_form}.
\end{proof}

Proposition~\ref{prop:search} reconciles the abstract with the conclusion. Part (ii) is
the abstract's claim: search alone suffices to eliminate the suboptimal points, because
what matters for the EGM step is continuity of the marginal value function, and search
delivers it. Part (iii) is the conclusion's claim: search alone does \emph{not} deliver
uniform smoothness, because differentiability of the marginal value function is a
strictly stronger requirement, and search does not deliver that. These are two
distinct properties and the paper should say so explicitly; as currently drafted, the
abstract and the conclusion appear to contradict each other.

\begin{proposition}[One-sidedness]
\label{prop:onesided}
$(\Phi^{S})''\equiv0$ on $\{\delta<0\}$ and $\varepsilon^{\ast}\equiv0$ on
$\{\delta\le0\}$. Hence under directed search: (i) the measure of agents who adjust
while the adjust branch is strictly dominated is exactly zero; and (ii) there is
exactly zero smoothing on the keep side of $\mathcal{K}_t$, the second derivative of
$v_t$ jumping from $0$ to $\big(\partial_m\Delta_t\big)^{2}/S$ across the threshold.
\end{proposition}

Proposition~\ref{prop:onesided} is simultaneously the paper's best economic argument
and the formal explanation of its main numerical difficulty. Part (i) is the property
that distinguishes search from logit: under logit, $P^{\adj}(0)=\tfrac12$ and
probability mass is placed on the dominated branch arbitrarily far from indifference,
which is the mechanism behind the finding that a large share of adjusters adjust for
taste rather than fundamental reasons. Part (ii) is the exact reason for the asymmetric
smoothing visible in Fig.~\ref{fig:fig2}(b): the asymmetry is a property of the
functional form, and \emph{no} choice of $S$ removes it. The paper attributes the
asymmetry to an economic mechanism (agents with low resources do not search); that
account is correct but incomplete, since the deeper cause is the truncation at
$\varepsilon^{\ast}=0$.

\begin{proposition}[The trade-off is not a theorem]
\label{prop:tradeoff}
Within the family of Lemma~\ref{lem:conjugate}, one-sidedness requires only
$\sigma'(0)=0$ while $\Phi\in C^{2}$ requires $\sigma''(0^{+})=+\infty$; the two
conditions are compatible. In particular, for the power cost
\begin{align}
\sigma_\gamma(\varepsilon) \;=\; A\,\frac{\gamma}{\gamma+1}\,\varepsilon^{(\gamma+1)/\gamma},
\qquad \gamma\in\mathbb{N},\ \gamma\ge 1,\ A>0,
\label{eq:powercost}
\end{align}
one has $\sigma_\gamma'(\varepsilon)=A\varepsilon^{1/\gamma}$ with $\sigma_\gamma'(0)=0$
for every $\gamma$, and $\sigma_\gamma''(0^{+})=+\infty$ for every $\gamma>1$. By
Lemma~\ref{lem:conjugate}(iv), $\Phi\in C^{\gamma}\setminus C^{\gamma+1}$, so
$v_t\in C^{\gamma}$, $c_{t-1}\in C^{\gamma-1}$, and marginal propensities to consume
are continuous for every $\gamma\ge 2$. The log-type
cost~\eqref{eq:search_cost_basic} satisfies $\sigma'(\varepsilon)\sim S\varepsilon$ and is
therefore the case $\gamma=1$: the \emph{least smooth} member of the one-sided family.
\end{proposition}

\begin{proof}
Immediate from Lemma~\ref{lem:conjugate}(ii)--(iv) and
$-S\log(1-\varepsilon)=S\varepsilon+O(\varepsilon^{2})$.
\end{proof}

Proposition~\ref{prop:tradeoff} is, I think, the most useful result in this appendix.
It shows that the apparent conflict between the economic advantage of search
(one-sidedness, hence no dominated-branch adjusters, hence an intact Euler-equation
interpretation away from the thresholds) and its numerical disadvantage (only $C^{1}$,
hence jumping marginal propensities, hence the need for taste shocks to obtain uniform
smoothness) is not a genuine trade-off. It is an artifact of a cost function whose
marginal cost vanishes linearly at zero. Adopting~\eqref{eq:powercost} with
$\gamma\ge2$ delivers both properties at once, and removes the need to combine search
with taste shocks --- which in turn removes the tension between
Section~\ref{sec:RetirementModel}, where the two are combined, and
Section~\ref{subsec:ModelBasicFull}, where combining them is rejected as
double-smoothing. The economic content of $\gamma$ is the elasticity of the search
hazard with respect to the value gap near the threshold, which is in principle
estimable from the same duration data used to discipline $S$.

%------------------------------------------------------------
\subsection{Uniformity in the horizon}
\label{subsec:smooth_uniform}
%------------------------------------------------------------

The results above are one-step statements. Because the paper's claim concerns the
accumulation of non-smoothness across backward iterations, bounds uniform in $t$ are
required. The following a priori estimate supplies them.

\begin{proposition}[Uniform curvature bound and explicit threshold]
\label{prop:uniform}
Let $\Phi$ satisfy $\|\Phi''\|_{\infty}\le C_\Phi<\infty$, so
$C_\Phi=1/S$ under directed search (Lemma~\ref{lem:searchclosed}) and
$C_\Phi=1/(4\sigma_\varepsilon)$ under EV Type~I taste shocks
(Proposition~\ref{prop:evtaste}). Suppose
Assumptions~\ref{as:u}--\ref{as:impatience} hold and $\Xi\equiv\sup_t\Xi_t<\infty$.
Then
\begin{align}
\big\|v_{mm,t}\big\|_{\infty}
&\le \beta\mathcal{R}^{2}\big\|v_{mm,t+1}\big\|_{\infty} + C_\Phi\,\Xi^{2},
\label{eq:apriori_step}\\[2pt]
\sup_{t\le T}\big\|v_{mm,t}\big\|_{\infty}
&\le \max\left\{\big\|v_{mm,T}\big\|_{\infty},\;
\frac{C_\Phi\,\Xi^{2}}{1-\beta\mathcal{R}^{2}}\right\}.
\label{eq:apriori_bound}
\end{align}
Consequently, by Lemma~\ref{lem:egm}(ii), the EGM step is clean at every $t$ whenever
\begin{align}
S \;>\; S^{\ast} \;\equiv\;
\frac{\beta\mathcal{R}^{2}\,\Xi^{2}}{\big(1-\beta\mathcal{R}^{2}\big)\,\inf\big|u''\big|},
\qquad\text{respectively}\qquad
\sigma_\varepsilon \;>\; \sigma_\varepsilon^{\ast} \;=\; S^{\ast}/4 .
\label{eq:threshold}
\end{align}
\end{proposition}

\begin{proof}
Differentiating the envelope identity $v_{m,t}^{j}=u'(c_t^{j})=\beta\mathcal{R}\mathcal{E}_t$
in $m$ gives
$v_{mm,t}^{j}=\beta\mathcal{R}^{2}\,\E_t[v_{mm,t+1}]\cdot\partial a_t/\partial m$, and
$\partial a_t/\partial m=1-\partial c_t/\partial m\in(0,1)$ by
Assumption~\ref{as:interior}, so
$\|v_{mm,t}^{j}\|_{\infty}\le\beta\mathcal{R}^{2}\|v_{mm,t+1}\|_{\infty}$. Differentiating
\eqref{eq:reduced_form} twice,
$v_{mm,t}=(1-\Phi')v_{mm,t}^{0}+\Phi'v_{mm,t}^{1}+\Phi''(\Delta_t)(\partial_m\Delta_t)^{2}$
with $\Phi'\in[0,1]$, which gives~\eqref{eq:apriori_step}. Iterating the affine
recursion under Assumption~\ref{as:impatience} and summing the geometric series
gives~\eqref{eq:apriori_bound}. Substituting into~\eqref{eq:egm_condition} and solving
for the smoothing parameter gives~\eqref{eq:threshold}.
\end{proof}

Two honest caveats. First, $\Xi<\infty$ uniformly in $t$ is an assumption, provable
only on a compact region of the state space; it should be verified numerically.
Second, as noted after Lemma~\ref{lem:egm}, $\inf|u''|$ vanishes at high wealth under
CRRA felicity, so~\eqref{eq:threshold} is informative only when stated regionally or in
consumption-equivalent units. Making the constants in~\eqref{eq:threshold}
non-vacuous is the substantive remaining work; the structure of the argument is
otherwise complete.

%------------------------------------------------------------
\subsection{The retirement model has a nested aggregator}
\label{subsec:smooth_retirement}
%------------------------------------------------------------

The retirement model of Section~\ref{sec:RetirementModel} does not fit the single-layer
template of Section~\ref{subsec:smooth_framework}, and the difference changes what can
be claimed. Collecting terms in the worker's problem, and using
$P_1(\varepsilon)=(1-\varepsilon)p+\varepsilon$ and $P_2(\varepsilon)=(1-\varepsilon)(1-p)$,
\begin{align}
\beta\big[P_1(\varepsilon)\goth{V}_t^{1}+P_2(\varepsilon)\goth{V}_t^{2}\big]
= \beta\big[p\,\goth{V}_t^{1}+(1-p)\,\goth{V}_t^{2}\big]
+ \varepsilon\cdot\beta(1-p)\big[\goth{V}_t^{1}-\goth{V}_t^{2}\big],
\end{align}
so that maximising over effort yields the reduced form
\begin{align}
\goth{V}_t(A_t)
&= \log c_t - \delta
+ \beta\big[p\,\goth{V}_t^{1}(A_t)+(1-p)\,\goth{V}_t^{2}(A_t)\big]
+ \Phi^{S}_{\sigma_s}\!\Big(\beta(1-p)\,\Delta_t^{\mathrm{opt}}(A_t)\Big),
\label{eq:ret_reduced}\\
\Delta_t^{\mathrm{opt}}(A_t)
&\equiv \goth{V}_t^{1}(A_t)-\goth{V}_t^{2}(A_t)
= \E_{y,\xi}\Big[\big(V_{t+1}^{\Ret}(M_{t+1})-\xi_{t+1}V_{t+1}^{W}(M_{t+1})\big)^{+}\Big],
\label{eq:ret_optiongap}
\end{align}
where $\Phi^{S}_{\sigma_s}$ is the conjugate of $\theta$ from
Lemma~\ref{lem:searchclosed} with $S$ replaced by $\sigma_s$.

\begin{proposition}[Search and taste shocks act on different margins]
\label{prop:retirement}
In the retirement model:
\begin{enumerate}
    \item Search does \emph{not} smooth the retire-versus-work maximum. That maximum
    appears inside~\eqref{eq:ret_optiongap} as the positive part $(\cdot)^{+}$, and is
    smoothed only by mollification over $(y_{t+1},\xi_{t+1})$, i.e.\ by the income and
    taste shocks, via Lemma~\ref{lem:mollify}.
    \item Search smooths a different margin, namely the arrival of the retirement
    option, by applying $\Phi^{S}_{\sigma_s}$ to the option-value gap
    $\Delta_t^{\mathrm{opt}}$.
    \item With exact integration over $(y,\xi)$, $\Delta_t^{\mathrm{opt}}\in C^{\infty}$
    by Lemma~\ref{lem:mollify}(i), and $\Phi^{S}_{\sigma_s}$ then reduces
    $\goth{V}_t$ to exactly $C^{1}$: \emph{search is the binding constraint on
    smoothness}. With $K$-node discretisation of $\xi$ the ranking reverses, by
    Proposition~\ref{prop:lognormaltaste}.
    \item Without taste or income uncertainty ($J=K=1$),
    $\Delta_t^{\mathrm{opt}}(A)=\big(V^{\Ret}-V^{W}\big)^{+}$ vanishes identically to the
    left of the retirement threshold, so $\varepsilon^{\ast}\equiv0$ there and
    $\goth{V}_t$ coincides with the no-search solution; to the right,
    $\Delta_t^{\mathrm{opt}}\sim\kappa\,(m-m^{\ast})^{+}$ and
    $\Phi^{S}_{\sigma_s}$ contributes
    $\big(\beta(1-p)\kappa\big)^{2}\big((m-m^{\ast})^{+}\big)^{2}/(2\sigma_s)$, a purely
    one-sided quadratic.
\end{enumerate}
\end{proposition}

Proposition~\ref{prop:retirement}(iv) reproduces Fig.~\ref{fig:fig2}(b) analytically:
the smoothing is exactly quadratic on one side of the threshold and exactly absent on
the other, with a second-derivative jump of $(\beta(1-p)\kappa)^{2}/\sigma_s$. Part
(iii) explains why taste shocks are indispensable in this application but not in the
durable one, where by contrast $\Phi^{S}$ \emph{replaces} the keep--adjust maximum
directly and Proposition~\ref{prop:search} applies verbatim. The claim in
Section~\ref{sec:DurableModel} that the durable exercise demonstrates ``the same level
of smoothness'' from search therefore requires qualification: in the durable model
search substitutes for the hard maximum, whereas in the retirement model it does not.

%------------------------------------------------------------
\subsection{The durable model in two dimensions}
\label{subsec:smooth_durable}
%------------------------------------------------------------

Three items require attention beyond the one-dimensional case.

\paragraph{Composite state.}
With $\Delta_t(m,n)=v_t^{\adj}\big(m+(1-\tau)n\big)-v_t^{\keep}(m,n)$, the reduced
form~\eqref{eq:reduced_form} is unchanged and
Propositions~\ref{prop:kink},~\ref{prop:evtaste},~\ref{prop:search},~\ref{prop:onesided}
and~\ref{prop:tradeoff} transfer verbatim with $\nabla$ in place of $\partial_m$, given
Assumptions~\ref{as:transversal}--\ref{as:singlecross}. The one-dimensionality of
$v_t^{\adj}$ in $\ell_t$ should be exploited: the adjuster's conditional value is a
function of a single scalar, so all statements about $\Phi$ apply along the
$\ell$-direction.

\paragraph{The adjuster's own non-convexity is a separate matter.}
Lemma~\ref{lem:egm} must be replaced by a Jacobian condition. Along the adjuster's
first-order locus, with $c(d,a)$ from $u_c=\mathfrak{v}'_c$ and $d$ solving
$u_d-\mathfrak{v}'_c+\mathfrak{v}'_d=0$, the endogenous liquidity map is
$\ell(d,a)=c(d,a)+d+a$, and refinement by FUES is required exactly when
$\partial\ell/\partial d$ changes sign. A sufficient condition for a clean adjuster
step is negative definiteness of
\begin{align}
\nabla^{2}_{(c,d)}u \;+\; \beta\,\nabla^{2}\,\E_t\big[v_{t+1}\big]
\label{eq:adj_jacobian}
\end{align}
on the relevant region, which by Proposition~\ref{prop:uniform} holds once
$\|\nabla^{2}v_{t+1}\|$ is small relative to $\big|\nabla^{2}u\big|$. This is where the
claim in the abstract and conclusion must be weakened: smoothing the keep--adjust
aggregator bounds $\nabla^{2}v_{t+1}$ and therefore makes the adjuster's program
locally concave \emph{once $S$ exceeds a threshold}, so ``the upper-envelope step
becomes unnecessary'' is provable, but only as a conditional statement with an explicit
threshold. It is not unconditional, and Section~\ref{sec:DurableModel} already observes
that residual discontinuities survive in the adjuster policy.

\paragraph{Concavity is not restored.}

\begin{proposition}[No smoother restores concavity]
\label{prop:noconcavity}
Let $\Phi$ be convex, nondecreasing and non-affine, as is the case for both
$\Phi^{\sigma}$ and $\Phi^{S}$. Then $v_t=v_t^{0}+\Phi(\Delta_t)$ is not concave in
general even when $v_t^{0}$ and $v_t^{1}$ are concave, since $\Phi$ composed with the
difference of two concave functions need not be concave. Concavity is recovered only in
the degenerate limits $\sigma_\varepsilon\to\infty$ or $S\to\infty$, in which the
discrete choice disappears.
\end{proposition}

Proposition~\ref{prop:noconcavity} should be stated explicitly rather than left open,
because it is the first question a reader will ask. It implies that the Euler equation
remains merely necessary in principle under either smoother, and therefore that the
correct target of the analysis is the bounded-curvature criterion
\eqref{eq:egm_condition} rather than sufficiency of the Euler equation. Proving the
weaker but operative result is more informative than leaving the concavity question
unaddressed.

%------------------------------------------------------------
\subsection{Summary of smoothness classes}
\label{subsec:smooth_summary}
%------------------------------------------------------------

Table~\ref{tab:smoothclasses} collects the results. Throughout, ``policy'' refers to
the \emph{conditional} policy $c_{t-1}$ of Remark~\ref{rem:twodiscont}(ii), income
integration is exact, and $\gamma\in\mathbb{N}$ indexes the power
cost~\eqref{eq:powercost}.

\begin{table}[htbp]
\centering
\caption{Smoothness class delivered by each aggregator. ``Clean EGM'' refers to
strict monotonicity of the endogenous grid map; ``bias'' is the per-period sup-norm
deviation of the aggregator from the hard maximum.}
\label{tab:smoothclasses}
\begin{tabular}{lccccc}
\toprule
Aggregator & $\Phi$ class & $v_t$ & $v_{m,t}$ & $c_{t-1}$ & MPC \\
\midrule
Hard max & $C^{0}\setminus C^{1}$ & $C^{0}$ & discontinuous & discontinuous & unbounded \\
EV Type~I, $\sigma_\varepsilon>0$ & analytic & $C^{\infty}$ & $C^{\infty}$ & $C^{\infty}$ & continuous \\
Log-normal, $K$ nodes & --- & $C^{0}$ & discontinuous & discontinuous & unbounded \\
Log-normal, exact & --- & $C^{\infty}$ & $C^{\infty}$ & $C^{\infty}$ & continuous \\
Search, log cost & $C^{1}\setminus C^{2}$ & $C^{1}$ & $C^{0}$ & $C^{0}$ & jumps \\
Search, power cost $\gamma$ & $C^{\gamma}\setminus C^{\gamma+1}$ & $C^{\gamma}$ & $C^{\gamma-1}$ & $C^{\gamma-1}$ & continuous if $\gamma\ge2$ \\
\midrule
& $\|\Phi''\|_\infty$ & bias & clean EGM & one-sided & \\
\midrule
Hard max & $+\infty$ & $0$ & never & yes & \\
EV Type~I, $\sigma_\varepsilon>0$ & $1/(4\sigma_\varepsilon)$ & $\le\sigma_\varepsilon\log 2$ & if $\sigma_\varepsilon>\sigma_\varepsilon^{\ast}$ & no & \\
Search, log cost & $1/S$ & $\le S$ & if $S>S^{\ast}$ & yes & \\
Search, power cost $\gamma$ & $+\infty$ & $O(A)$ & see text & yes & \\
\bottomrule
\end{tabular}
\end{table}

The last row requires a comment. For $\gamma>1$ the power cost has
$\sigma''(0^{+})=+\infty$ and therefore $\Phi''$ unbounded at the origin, so the
uniform bound~\eqref{eq:apriori_bound} does not apply directly; the correct statement
replaces the sup-norm bound on $\Phi''$ by a bound on
$\Phi''(\Delta_t)(\partial_m\Delta_t)^{2}$ over the relevant range of $\Delta_t$, which
is finite away from an arbitrarily small neighbourhood of $\mathcal{K}_t$ and
integrable across it. Establishing the analogue of
Proposition~\ref{prop:uniform} for $\gamma>1$ is the one genuinely open technical item
in this appendix.

%------------------------------------------------------------
\subsection{Numerical tests implied by the theory}
\label{subsec:smooth_tests}
%------------------------------------------------------------

The propositions above make sharp, falsifiable predictions that can replace the visual
evidence currently offered in Sections~\ref{sec:RetirementModel}
and~\ref{sec:DurableModel}. Each is computable with the existing solver.

\begin{enumerate}
    \item \textbf{Conservation of total jump} (Lemma~\ref{lem:mollify}(iii)). The total
    variation of $c_t'$ should be invariant in the number of income nodes $J$ while the
    largest individual jump falls at rate $O(1/J)$. This is the correct diagnostic for
    Section~\ref{sec:incdisc}, and it is strictly more informative than the maximum
    absolute deviation reported in Table~\ref{tab:incdisc}, which by
    Proposition~\ref{prop:exactincome} is dominated by threshold placement.

    \item \textbf{The monotonicity threshold} (Lemma~\ref{lem:egm},
    Proposition~\ref{prop:uniform}). Report $\min_a m_t'(a)$ as a function of $S$ and of
    $\sigma_\varepsilon$, and locate the value at which it crosses zero. Compare the
    empirical threshold with the analytic $S^{\ast}$ of~\eqref{eq:threshold}. This
    replaces ``the figures look smooth'' with a measured statement, and identifies the
    friction required to justify dropping the envelope step.

    \item \textbf{The second-difference signature} (Propositions~\ref{prop:evtaste}
    and~\ref{prop:search}). Across $\mathcal{K}_t$, the second difference of $v_t$
    should be continuous under EV Type~I taste shocks and should jump by exactly
    $\big(\partial_m\Delta_t\big)^{2}/S$ under directed search. This is a quantitative
    validation of the theory against the solver, and it distinguishes the two
    mechanisms unambiguously.

    \item \textbf{Uniformity in the horizon} (Proposition~\ref{prop:uniform}). Report
    $\|v_{mm,t}\|_{\infty}$ and $|\mathcal{N}_t|$ by $t$ for a sweep over $S$, and check
    that both stabilise rather than growing. This is the direct test of the accumulation
    claim, and it requires a horizon considerably longer than the five periods currently
    used for the durable figures.

    \item \textbf{The share of dominated-branch adjusters}
    (Proposition~\ref{prop:onesided}). At the $(S,\sigma_\xi)$ or
    $(S,\sigma_\varepsilon)$ combination that first delivers a clean EGM step, report the
    fraction of adjustments occurring where the adjust branch is strictly dominated.
    Under pure search this is exactly zero; under any admixture of taste shocks it is
    positive. This is the decisive quantitative test of the paper's economic argument
    against a taste-shock formulation and it is currently absent.
\end{enumerate}

\begin{remark}[Implementation prerequisite]
\label{rem:implementation}
The results of Section~\ref{subsec:smooth_search} concern the exact
conjugate~\eqref{eq:PhiS}. The accompanying implementation selects search intensity by
maximisation over a finite grid on $[0,1]$ rather than by the closed
form~\eqref{eq:eps_star_basic}. A grid-selected $\varepsilon$ makes the realised
aggregator piecewise affine and reintroduces kinks of order the grid spacing, which
would contaminate test~3 above in particular. The closed form should be used, or the
two should be verified to agree well below the magnitudes being measured.
\end{remark}

%------------------------------------------------------------
\subsection{Open items}
\label{subsec:smooth_open}
%------------------------------------------------------------

Ranked by difficulty. (i) Uniform boundedness of $\Xi_t$ and of $|\mathcal{N}_t|$ across
the backward recursion, and non-vacuous constants in~\eqref{eq:threshold}: the one-step
and invariant-set arguments are complete, but the constants should be calibrated
numerically rather than bounded analytically. (ii) The analogue of
Proposition~\ref{prop:uniform} for the power cost~\eqref{eq:powercost} with $\gamma>1$,
where $\Phi''$ is unbounded at the origin. (iii)
Assumption~\ref{as:singlecross} in the durable model, which should be verified on the
solved model rather than proved from primitives. Everything else ---
Lemmas~\ref{lem:transfer}--\ref{lem:searchclosed} and
Propositions~\ref{prop:kink}--\ref{prop:noconcavity} --- is self-contained and proved
above.

%============================================================
\subsection{Glossary of mathematical properties}
\label{subsec:smooth_glossary}
%============================================================

This subsection collects, for reference, precise statements of every property invoked
above, together with a pointer to where each is used. Throughout, $U\subseteq\mathbb{R}^{n}$
is open, $f:U\to\mathbb{R}$, and $\lvert\cdot\rvert$ is the Euclidean norm.

%------------------------------------------------------------
\subsubsection{Smoothness and singularity classes}
%------------------------------------------------------------

\begin{description}

\item[\textnormal{\textbf{Differentiability class $C^{k}$.}}]
$f\in C^{k}(U)$ if all partial derivatives of $f$ of order $\le k$ exist and are
continuous on $U$. By convention $C^{0}(U)$ is the continuous functions and
$C^{\infty}(U)=\bigcap_{k\ge0}C^{k}(U)$. We write $f\in C^{k}\setminus C^{k+1}$ to mean
that $f$ is \emph{exactly} of class $C^{k}$, i.e.\ $f\in C^{k}$ but some $k+1$-st partial
derivative fails to exist or fails to be continuous somewhere in $U$. All statements of
the form ``$v_t\in C^{1}\setminus C^{2}$'' (Proposition~\ref{prop:search}) are exactness
claims in this sense, and are what distinguishes the mechanisms from one another.

\item[\textnormal{\textbf{Real analyticity.}}]
$f$ is real analytic on $U$ if in a neighbourhood of each point of $U$ it equals a
convergent power series. Real analyticity is strictly stronger than $C^{\infty}$ and is
the class of the Extreme Value Type~I log-sum $\Phi^{\sigma}$
(Proposition~\ref{prop:evtaste}).

\item[\textnormal{\textbf{Lipschitz continuity.}}]
$f$ is Lipschitz on $U$ with modulus $L$ if $\lvert f(x)-f(y)\rvert\le L\lvert x-y\rvert$
for all $x,y\in U$. A pointwise maximum of finitely many $C^{1}$ functions with locally
bounded gradients is Lipschitz but generally not $C^{1}$; this is the exact sense in
which the hard-max value function of Proposition~\ref{prop:kink} is ``well behaved but
not smooth,'' and it is why the failure is invisible to a plot of $v_t$ itself and
visible only in $v_{m,t}$.

\item[\textnormal{\textbf{One-sided derivatives and the jump bracket.}}]
Let $\gamma:(-\epsilon,\epsilon)\to U$ be $C^{1}$ with $\gamma(0)=x^{\ast}$. The one-sided
directional derivatives of $f$ at $x^{\ast}$ along $\gamma$ are
\begin{align}
\partial^{\pm}f(x^{\ast};\gamma) \;\equiv\; \lim_{s\to 0^{\pm}}\frac{d}{ds}f(\gamma(s)),
\end{align}
when they exist, and the \emph{jump bracket} is
$\llbracket f'\rrbracket(x^{\ast};\gamma)\equiv\partial^{+}f-\partial^{-}f$. The jump
bracket of the second derivative, $\llbracket f''\rrbracket$, is defined analogously.
Equation~\eqref{eq:kink_jump} is a statement about $\llbracket v_t'\rrbracket$, and
Lemma~\ref{lem:searchclosed} is a statement about $\llbracket (\Phi^{S})''\rrbracket=1/S$.

\item[\textnormal{\textbf{Kink; convex and concave kinks.}}]
$f$ has a \emph{kink} at $x^{\ast}$ along $\gamma$ if it is continuous there and
$\llbracket f'\rrbracket\neq0$. The kink is \emph{convex} if
$\llbracket f'\rrbracket>0$ and \emph{concave} if $\llbracket f'\rrbracket<0$. The
distinction is not cosmetic: only a convex kink contributes a \emph{positive} singular
part to $f''$, and only a positive contribution
violates~\eqref{eq:egm_condition}. Proposition~\ref{prop:kink} shows that a hard
maximum always produces a convex kink, which is why the failure of the EGM step under a
hard maximum is systematic rather than parameter dependent.

\item[\textnormal{\textbf{Second-order kink.}}]
$f$ has a second-order kink at $x^{\ast}$ if $f\in C^{1}$ near $x^{\ast}$ but
$\llbracket f''\rrbracket\neq0$. This is the exact defect of directed search under the
log-type cost (Proposition~\ref{prop:search}(i)): the level and the gradient are
continuous, the curvature is not.

\item[\textnormal{\textbf{Singular set.}}]
$\mathcal{N}(f)\equiv\{x\in U: f\text{ is not }C^{1}\text{ at }x\}$, and
$\mathcal{N}_t\equiv\mathcal{N}(v_t)$. Proposition~\ref{prop:accumulate} is a statement
about the growth of $\mathcal{N}_t$ as $t$ recedes from $T$, and is the formal content
of the claim that kinks propagate backwards.

\item[\textnormal{\textbf{Piecewise class $\mathcal{PC}^{k}(\mathcal{N})$.}}]
For a closed set $\mathcal{N}\subset U$ of Lebesgue measure zero,
$f\in\mathcal{PC}^{k}(\mathcal{N})$ if $f\in C^{k}(U\setminus\mathcal{N})$ and all
partial derivatives of order $\le k$ admit finite one-sided limits at every point of
$\mathcal{N}$. This scale is needed because the hard-max case produces objects that are
not in any $C^{k}$ with $k\ge1$, so the $C^{k}$ scale alone cannot distinguish
``one kink'' from ``a thousand kinks'' --- a distinction on which
Proposition~\ref{prop:accumulate} turns.

\item[\textnormal{\textbf{Distributional derivative and Dirac mass.}}]
If $f$ is $C^{1}$ away from $x^{\ast}$ and $f'$ has jump $\kappa$ at $x^{\ast}$, then the
second derivative of $f$ in the sense of distributions is the measure
\begin{align}
f'' \;=\; \{f''\} \;+\; \kappa\,\delta_{x^{\ast}},
\end{align}
where $\{f''\}$ is the classical part and $\delta_{x^{\ast}}$ is the Dirac measure at
$x^{\ast}$. This is the precise sense in which $\E_t[v_{mm,t+1}]$ is unbounded at a kink
and in which condition~\eqref{eq:egm_condition} is ``violated infinitely badly'' under a
hard maximum with quadrature (Proposition~\ref{prop:quadfail}). Under exact income
integration the Dirac is convolved into the bounded spike of
Proposition~\ref{prop:exactincome}.

\end{description}

%------------------------------------------------------------
\subsubsection{Geometry of the indifference set}
%------------------------------------------------------------

\begin{description}

\item[\textnormal{\textbf{Regular value and transversality.}}]
Let $\Delta\in C^{1}(U)$. The value $0$ is a \emph{regular value} of $\Delta$ if
$\nabla\Delta(x)\neq0$ for every $x\in\Delta^{-1}(0)$. A $C^{1}$ curve $\gamma$
\emph{crosses} $\mathcal{K}=\Delta^{-1}(0)$ \emph{transversally} at $\gamma(0)=x^{\ast}$
if $\tfrac{d}{ds}\Delta(\gamma(s))\big|_{s=0}\neq0$, i.e.\ if $\dot\gamma(0)$ is not
tangent to $\mathcal{K}$. Assumption~\ref{as:transversal} is exactly the requirement
that $0$ be a regular value of $\Delta_t$ for every $t$.

\item[\textnormal{\textbf{Codimension-one $C^{1}$ submanifold.}}]
$\mathcal{K}\subset\mathbb{R}^{n}$ is a $C^{1}$ embedded submanifold of codimension one
if each of its points has a neighbourhood in which $\mathcal{K}$ is the graph of a
$C^{1}$ function of $n-1$ of the coordinates. By the regular value theorem (a corollary
of the implicit function theorem), if $0$ is a regular value of $\Delta$ then
$\Delta^{-1}(0)$ is such a submanifold. In the retirement model this reduces to a
finite set of isolated threshold points $M^{\ast}$; in the durable model it is the
$(S,s)$ boundary curve in $(m,n)$ space, and the theorem is what guarantees that the
boundary can be written locally as $m=\bar m(n)$ with $\bar m\in C^{1}$.

\item[\textnormal{\textbf{Jump magnitude $\kappa_t$ and its supremum $\Xi_t$.}}]
$\kappa_t=\inf_{\mathcal{K}_t}\lvert\nabla\Delta_t\rvert$ and
$\Xi_t=\sup_{\mathcal{K}_t}\lvert\nabla\Delta_t\rvert$. By
Proposition~\ref{prop:kink} the marginal value function jumps by exactly
$\lvert\nabla\Delta_t\rvert$ across $\mathcal{K}_t$ in the crossing direction, so
$\kappa_t>0$ is what makes the negative result non-vacuous and $\Xi_t<\infty$ is what
makes the positive results in Propositions~\ref{prop:evtaste},~\ref{prop:search}
and~\ref{prop:uniform} quantitative. In the retirement model
$\kappa_t=\lvert\lambda_t^{W}-\lambda_t^{\Ret}\rvert$ evaluated at the threshold.

\item[\textnormal{\textbf{Single crossing.}}]
$\Delta_t(\cdot,n)$ satisfies single crossing in $m$ if it changes sign at most once and,
where it does, from negative to positive; equivalently, $\{m:\Delta_t(m,n)>0\}$ is an
interval unbounded above. Single crossing at a regular value implies
$\partial_m\Delta_t>0$ at the zero, which is the form used in
Assumption~\ref{as:singlecross}. Economically this is what delivers an $(S,s)$
\emph{band}: adjustment occurs on an interval of $m$ rather than on an arbitrary set.

\end{description}

%------------------------------------------------------------
\subsubsection{Convexity, duality, and control costs}
%------------------------------------------------------------

\begin{description}

\item[\textnormal{\textbf{Strict concavity; concave program.}}]
$f$ is strictly concave on a convex $U$ if
$f(\lambda x+(1-\lambda)y)>\lambda f(x)+(1-\lambda)f(y)$ for all $x\neq y$ and
$\lambda\in(0,1)$. A \emph{concave program} maximises a concave objective over a convex
feasible set; in such a program the Karush--Kuhn--Tucker conditions are necessary
\emph{and} sufficient, and the maximiser is unique if the objective is strictly concave.
The entire difficulty addressed in this paper is that the conditional problems are
concave programs only if the continuation value is concave, and by
Proposition~\ref{prop:noconcavity} no smoother restores that concavity, which is why
the operative target is~\eqref{eq:egm_condition} rather than sufficiency of the Euler
equation.

\item[\textnormal{\textbf{Inada conditions.}}]
$u'(0^{+})=+\infty$ and $u'(+\infty)=0$ in each argument. These guarantee that the
optimal continuous choice is interior, so that the kinks studied here are attributable
to the discrete choice rather than to a corner (Assumption~\ref{as:interior}).

\item[\textnormal{\textbf{Translation equivariance.}}]
$\Lambda:\mathbb{R}^{2}\to\mathbb{R}$ is translation equivariant if
$\Lambda(v^{0}+k,v^{1}+k)=\Lambda(v^{0},v^{1})+k$ for every $k\in\mathbb{R}$. All three
aggregators in~\eqref{eq:three_Phi} have this property, and it is exactly what permits
the one-dimensional reduced form~\eqref{eq:reduced_form}. Without it, regularity would
have to be analysed in two arguments rather than one and the appendix would not
simplify.

\item[\textnormal{\textbf{Convex conjugate (Legendre--Fenchel transform).}}]
For $\sigma:[0,1]\to\mathbb{R}$, the convex conjugate is
\begin{align}
\sigma^{\ast}(\delta) \;\equiv\; \sup_{\varepsilon\in[0,1]}\big\{\varepsilon\delta-\sigma(\varepsilon)\big\}.
\end{align}
$\sigma^{\ast}$ is always convex and lower semicontinuous. If $\sigma$ is strictly
convex and $C^{2}$ on $(0,1)$ then $\sigma^{\ast}$ is differentiable with
$(\sigma^{\ast})'=(\sigma')^{-1}$ and
$(\sigma^{\ast})''(\delta)=1/\sigma''\big((\sigma')^{-1}(\delta)\big)$
\citep{Rockafellar1970-cv}. Lemma~\ref{lem:conjugate} is a direct application, and the
inversion formula for the second derivative is the entire content of
Proposition~\ref{prop:tradeoff}: \emph{curvature of the smoother is the reciprocal of
curvature of the cost}.

\item[\textnormal{\textbf{Danskin's theorem (envelope theorem for maxima).}}]
If $g(\varepsilon;\delta)$ is $C^{1}$ in $\delta$, concave in $\varepsilon$, and the
maximiser $\varepsilon^{\ast}(\delta)$ is unique, then
$\tfrac{d}{d\delta}\max_{\varepsilon}g=\partial_{\delta}g\big(\varepsilon^{\ast}(\delta);\delta\big)$:
the indirect effect through $\varepsilon^{\ast}$ vanishes. This justifies
$\Phi'=\varepsilon^{\ast}$ in Lemma~\ref{lem:conjugate}(i) and the combined
marginal-value formulas~\eqref{eq:search_vm}--\eqref{eq:search_vn} of
Section~\ref{subsec:appdx_search_basic}.

\item[\textnormal{\textbf{Control cost; perturbed utility.}}]
A \emph{control cost} is a strictly convex function $\sigma$ on the simplex of
randomisations over the discrete alternatives, subtracted from the expected payoff, so
that the agent solves $\max_{\varepsilon}\{\varepsilon\delta-\sigma(\varepsilon)\}$. The
resulting choice probability is $\varepsilon^{\ast}=(\sigma')^{-1}(\delta)$. When
$\sigma$ is the negative Shannon entropy the solution is logit, which is the sense in
which Extreme Value Type~I taste shocks and directed search are members of one family
\citep{Mattsson2002-pc,Hofbauer2002-gc,Fudenberg2015-sv}. Additive random-utility and
additive perturbed-utility representations of stochastic choice are observationally
equivalent, which is why the ``rationalised versus behavioural'' contrast cannot by
itself distinguish the two mechanisms; see the discussion after
Lemma~\ref{lem:conjugate}.

\item[\textnormal{\textbf{One-sidedness of a smoother.}}]
$\Phi$ is \emph{one-sided} if $\Phi'(\delta)=0$ for every $\delta\le0$, i.e.\ zero weight
is ever placed on a strictly dominated branch. By Lemma~\ref{lem:conjugate}(ii) this
holds if and only if the control cost has zero marginal cost at zero intensity,
$\sigma'(0)=0$. Directed search is one-sided; logit is not
(Proposition~\ref{prop:onesided}).

\item[\textnormal{\textbf{Inclusive value.}}]
The certainty-equivalent value of a discrete choice before the shocks are realised,
$\E\big[\max_j\{v^{j}+\text{shock}_j\}\big]$. For Extreme Value Type~I shocks this is
the log-sum~\eqref{eq:taste_logsum}. Note $\Phi^{\sigma}\ge\Phi^{\max}$ (the inclusive
value exceeds the maximum, an option value) whereas $\Phi^{S}\le\Phi^{\max}$ (search
absorbs resources); see the remark following Lemma~\ref{lem:searchclosed}.

\end{description}

%------------------------------------------------------------
\subsubsection{Integration and approximation}
%------------------------------------------------------------

\begin{description}

\item[\textnormal{\textbf{Mollification (convolution smoothing).}}]
Given a density $f$ and scale $y>0$, the mollification of $g$ is
$G(b)=\int g(b+y\eta)f(\eta)\,d\eta=y^{-1}\int g(z)f\big((z-b)/y\big)dz$. The essential
property, and the content of Lemma~\ref{lem:mollify}(i), is that differentiation in $b$
falls on $f$ rather than on $g$, so $G$ inherits the smoothness of the \emph{smoother}
of its two factors. Exact integration of a $C^{\infty}$ density therefore produces a
$C^{\infty}$ function no matter how badly behaved $g$ is.

\item[\textnormal{\textbf{Quadrature rule; equiprobable discretisation.}}]
A quadrature rule replaces $\int h\,dF$ by $\sum_{j\le J}\omega_j h(\eta_j)$ with
$\omega_j>0$ and $\sum_j\omega_j=1$; an equiprobable discretisation sets
$\omega_j=1/J$ with $\eta_j$ the conditional means of $J$ equiprobable bins. A
quadrature rule is a finite convex combination of \emph{evaluations} and therefore, by
Lemma~\ref{lem:mollify}(ii), has \emph{no} mollifying effect: it inherits every
singularity of $g$, replicated once per node. The distinction between this entry and the
previous one is the rigorous content of the phantom-kink discussion in
Section~\ref{sec:incdisc}.

\item[\textnormal{\textbf{Total variation.}}]
For $h$ on an interval $[\underline x,\bar x]$,
$\mathrm{TV}(h)=\sup\sum_{i}\lvert h(x_{i+1})-h(x_i)\rvert$ over all finite partitions;
for piecewise $C^{1}$ $h$ this equals $\int\lvert h'\rvert$ plus the sum of the absolute
jumps. Lemma~\ref{lem:mollify}(iii) states that $\mathrm{TV}$ of the derivative is
conserved under refinement of a quadrature rule, which makes $\mathrm{TV}(c_t')$ the
correct diagnostic for numerically induced roughness --- strictly more informative than
the maximum absolute deviation reported in Table~\ref{tab:incdisc}.

\end{description}

%------------------------------------------------------------
\subsubsection{Properties of the EGM step}
%------------------------------------------------------------

\begin{description}

\item[\textnormal{\textbf{Endogenous grid map.}}]
$m_t(a)=a+c_t(a)$, where $c_t(a)=(u')^{-1}\big(\beta\mathcal{R}\mathcal{E}_t(a)\big)$ is
obtained by inverting the Euler equation at an exogenous end-of-period asset value $a$.
The method is well posed precisely when this map can be inverted to recover $a$, and
hence $c_t$, as a function of $m_t$.

\item[\textnormal{\textbf{Clean EGM step; suboptimal point.}}]
The EGM step at $t$ is \emph{clean} if $a\mapsto m_t(a)$ is continuous and strictly
increasing on the unconstrained region (Definition~\ref{def:clean}). A
\emph{suboptimal point} is a candidate pair $\big(m_t(a),c_t(a)\big)$ satisfying the
Euler equation but not attaining the maximum; such points exist exactly when $m_t$
fails to be injective, since then two or more values of $a$ deliver the same $m_t$ and
at most one can be optimal. Lemma~\ref{lem:egm} characterises cleanliness by the
curvature condition~\eqref{eq:egm_condition}, and this --- not differentiability --- is
what the abstract's claim about eliminating suboptimal points should be understood to
assert.

\item[\textnormal{\textbf{Upper envelope refinement.}}]
Any procedure that, given the candidate cloud $\{(m_t(a),c_t(a),v_t(a))\}$, discards the
dominated candidates and returns the upper envelope of the implied value
correspondence; DC-EGM \citep{Iskhakov2017-jb} and FUES
\citep{Dobrescu2022-ii} are instances. A clean EGM step is exactly the condition under
which this refinement is unnecessary.

\end{description}

%------------------------------------------------------------
\subsubsection{Dynamic-programming properties}
%------------------------------------------------------------

\begin{description}

\item[\textnormal{\textbf{Envelope bootstrap.}}]
The observation, formalised in Lemma~\ref{lem:transfer}(ii), that the envelope condition
$v_{m,t}^{j}=u'(c_t^{j})$ recovers one order of differentiability relative to the Euler
equation, so that conditional value functions \emph{preserve} their smoothness class
across a backward step. It follows that only the aggregator can degrade smoothness, and
this is what reduces the entire induction to a statement about $\Phi$ near the origin.

\item[\textnormal{\textbf{Invariant smoothness class.}}]
A class $\mathcal{C}$ is invariant for the backward operator if $v_{t+1}\in\mathcal{C}$
implies $v_{t}\in\mathcal{C}$. The three mechanisms have invariant classes $C^{0}$
(hard maximum, with a growing singular set), $C^{1}$ (directed search, log-type cost),
and $C^{\infty}$ (Extreme Value Type~I). Invariance is what upgrades a one-step
regularity result into a statement about accumulation over the horizon.

\item[\textnormal{\textbf{A priori estimate; invariant set.}}]
An inequality bounding a solution in terms of the data before the solution is known.
Proposition~\ref{prop:uniform} establishes one by showing that the backward operator
maps $\{v:\lVert v_{mm}\rVert_{\infty}\le B\}$ into itself for
$B=C_\Phi\Xi^{2}/(1-\beta\mathcal{R}^{2})$; the set is then invariant and the bound holds
uniformly in $t$.

\item[\textnormal{\textbf{Contraction modulus.}}]
The factor $\beta\mathcal{R}^{2}<1$ (Assumption~\ref{as:impatience}) by which curvature
is discounted per backward step, arising as $\beta\mathcal{R}$ from discounting the
marginal value and a further $\mathcal{R}$ from the chain rule through
$m_{t+1}=\mathcal{R}a_t+\theta_{t+1}$. It is what makes the geometric series
in~\eqref{eq:apriori_bound} converge, and hence what rules out unbounded accumulation of
curvature over long horizons.

\end{description}
